\section{A degree-separated nine-variable equation}
\label{sec:nine-degree-separated}

The primary equation has nine strictly positive unknowns, exact total degree
\(2\,676\,548\), and \(476\) live binary operations.  Its separately
translated nonnegative form has the same degree.  The complete division-free
DAGs, coefficient payloads, exact-degree certificates, and a standard-library
checker accompany this section.  Register names below abbreviate polynomials;
they introduce no additional unknowns.

\paragraph{Scope.}
The ordinary mathematical theorem is that ZFC proves existence of a zero in
the stated domain equivalent to existence of a finite syntactic ZFC
refutation.  Its negation is the requested consistency equivalence.
The original package31 compiler and the numerical Jones ratio theorems are
explicit imported dependencies.  This is not a proof of consistency, a new
end-to-end kernel formalization, an unrestricted-integer nine-variable claim,
or a global minimum claim.

\subsection{A smaller source by keeping the three degrees separate}
Use the original 31-positive-variable, 22-residual source
\(r_1,\ldots,r_{22}\), with its \emph{whole} original coefficient tuple.
The four payload bit lengths are
\(1064921,1070004,1064921,1064922\); the fourth port is already shifted.
Replace just
\[
 \rho_{14}=r_{14}+r_2,\qquad \rho_{15}=r_{15}+r_3.
\]
The inverse row operations recover the old system, so its zero set is
unchanged.  It has 89 weighted records: 7 constant, 39 linear, and 43 quadratic.

For positive source coordinates \(y_i\), write
\(U_y(t)=\sum_i y_i t^{w_i}\), with weights
\begin{align*}
 &(1,7,61,36,16,85,20,67,43,11,34,13,96,80,29,50,\\
 &\hspace{8mm}2,109,110,112,116,91,53,55,27,101,92,9,107,37,60).
\end{align*}
The maximum weight is \(W=116\).  For a monomial of degree \(d\), weight its
coefficient by 2 when \(d=0,1\), and by \(\prod_i\alpha_i!\) when \(d=2\).
With the fixed target assigned to its row, form
\[
 D_d(t)=\sum_{j,|\alpha|=d}c_{j,\alpha}
 t^{q_j-\sum_i\alpha_iw_i},\qquad
 \mathcal G(t,y)=D_2(t)U_y(t)^2+D_1(t)U_y(t)+D_0(t).
\]
The finite target certificate makes a degree-\(d\) term hit a marked target
only against its own monomial of that same degree.  Cross-degree label
collisions are harmless, because the three kernels are separate summands.
Checking all \(22\,544\) interactions proves the formal identities
\[
 [t^{q_j}]\mathcal G(t,y)=2\rho_j(y).
\]
The kernel degrees are \((1719,2027,1929)\), the leading coefficient of
\(D_2\) is \(+1\), and the convolution height is
\(K_{\rm code}=2161\).  The largest target is 2129.  All 89 kernel records,
coefficient recipes, and the row-target map are supplied as exact JSON.

\subsection{Positive words and unconditional source bounds}
Every weighted coefficient has absolute value below \(M=2^{2134927}\).
Set the fixed ordinary integer coefficients
\[
 a=4096M=2^{2134939},\qquad c=2^{17079502}.
\]
Use mixed coordinates \(b,g,h,s,w,i\ge1\) and \(\phi,j,n\ge0\).  Define
\begin{align*}
 B&=ab^2,&\ell&=\sum_i B^{w_i},& U_s&=\ell+g,\\
 G_s&=(D_2(B)U_s+D_1(B))U_s+D_0(B),&
 A_q&=\sum_j B^{q_j},\\
 N_0&=B^{117},& Z&=B^{2162},&
 M_c&=N_0-1-(b-1)\ell,\\
 J_c&=2G_s+A_q,&S_c&=g+N_0J_c,&
 T_c&=M_c+N_0(B-2)A_q,\\
 L&=2bB^{2279},&N&=cbB^{1201},&G_{\rm cap}&=32bB^{116},\\
 R&=L(L+S_c+T_c+1)-S_c-1.&&&
\end{align*}
Here \(L\) is the binary-block length, while \(N\) is the separate
square-root threshold for Jones's theorem.

Positivity must precede digit recovery.  For every \(g\ge1\),
\(U_s\ge B^{116}\), and geometric bounds give
\[
 D_2(B)>\tfrac12B^{1929},\quad
 |D_1(B)|<2MB^{2027},\quad |D_0(B)|<2MB^{1719}.
\]
Relative to \(B^{1929}U_s^2\), the negative contributions are bounded by
\(2M/B^{18}\) and \(2M/B^{442}\), each below \(1/8\).  Thus
\(G_s>\tfrac14B^{1929}U_s^2>0\).  Also \(0\le M_c<N_0\), and
\[
 R=L^2+(L-1)S_c+LT_c+L-1\ge L^2+L-1.
\]
Consequently \(R,N\ge8\), \(b\le R,N\), and \(G_{\rm cap}>0\)
on every outer assignment, without first assuming the final equation vanishes.

On the branch \(g<4bB^{116}\), one has \(U_s<6bB^{116}\), and
\[
 |G_s|<\frac{43Mb^2 B^{2162}}{B-1},\qquad
 2G_s<Z/32,\qquad A_q,(B-2)A_q<Z/64.
\]
Hence \(g<N_0\), \(J_c<3Z/64\), and
\[
 S_c+T_c<N_0Z=B^{2279}<L.
\]
These are numerical, prebinary bounds, not circular consequences of the carry
test.

If \(b\) is a power of two, so is \(B\).  The low-mask test
\(g\mathbin{\&}M_c=0\) says exactly
\(g=\sum_i z_iB^{w_i}\), \(0\le z_i<b\).  Then \(y_i=z_i+1\) are
positive and \(U_s=\sum_i y_iB^{w_i}\).  Each coefficient \(h_t\) of
\(\mathcal G(t,y)\), including unsupported collisions, satisfies
\[
 |h_t|<M\bigl((\textstyle\sum_i y_i)^2+\sum_i y_i+1\bigr)
 \le993Mb^2<B/4.
\]
The signed digits of \(J_c\) are \(2h_t+1\) at targets and \(2h_t\)
elsewhere.  Normalization has incoming and outgoing borrows in \(\{-1,0\}\).
At a target, the normalized digit avoids \(B-2\) exactly when \(h_t=0\):
the allowed digits are 0 or 1, while positive or negative \(h_t\) activate a
masked bit.  Therefore
\[
 S_c\mathbin{\&}T_c=0
 \quad\Longleftrightarrow\quad
 g\text{ encodes a positive common zero of the source.}
\]
Conversely, any source zero admits every sufficiently large power-of-two
\(b\); its digits are \(z_i=y_i-1\).  The all-one source tuple is not a
zero, because \(r_5(1,\ldots,1)=1\).  Thus \(g>0\), and
\(g<bB^{116}\), giving the stronger completeness bound
\(32g<G_{\rm cap}\).

\subsection{A calibrated two-block threshold}
For \(S,T\ge0\), \(S+T<L=2^u\), the integer
\[
 X=L^2+(S+T)L+(L-1-S)
\]
has nonoverlapping binary blocks.  Consequently
\[
 \operatorname{pop}(X)=u+\operatorname{pop}(T)+1
 -\bigl[\operatorname{pop}(S)+\operatorname{pop}(T)
 -\operatorname{pop}(S+T)\bigr].
\]
The bracket is nonnegative and vanishes exactly for carry-free addition.
Legendre's formula identifies \(\operatorname{pop}(X)\) with
\(v_2\binom{2X}{X}\).

Write \(b=2^k\), \(B=2^v\), \(v=2134939+2k\).  The unchanged low
mask has \(117v-31k\) set bits and the 22 target cells contribute
\(22(v-1)\).  Hence
\[
 \operatorname{pop}(T_c)=139v-31k-22,\qquad
 u=1+k+2279v.
\]
The exact threshold is
\[
 u+\operatorname{pop}(T_c)+1=2418v-30k-20
 =2(17079502+k+1201v)=2\log_2N.
\]
The factor \(b\) in \(N\) accommodates the even parameter 2418, without
restricting \(b\) to a sparse family.  We have proved
\[
 N^2\mid\binom{2R}{R}\quad\Longleftrightarrow\quad
 S_c\mathbin{\&}T_c=0.
\]

\subsection{Ratio registers and the positivity-first finalizer}
Define
\begin{align*}
 U&=N^2w,&Y&=N^2s,&M_0&=RY,\\
 P&=2M_0^2U,&K&=R+1+h(P-1),&A&=M_0(U+1),\\
 B_1&=2R+1,&W_1&=bw,&C&=B_1+W_1+\phi,\\
 J&=4A-5,&Q&=3W_1C-2(W_1^2-1),\\
 D&=(A^2-1)C^2+1,&k_0&=iJC,&E&=k_0DC,\\
 F&=(A^2-1)E^2+1,&G&=A+F(CD+1-A),\\
 H&=C+F(B_1+jC),&I&=(G^2-1)H^2+1,\\
 a_0&=(P^2-1)K^2+1,&a_1&=DFI.
\end{align*}
Since \(F\equiv1\pmod J\) and \(F\equiv1\pmod C\), the old
divisibility \(JF\mid QF+J(H_{\rm old}-C)\) is equivalent to
\(J\mid Q\) and \(F\mid H_{\rm old}-C\).  The exact witness substitution
\[
 H=C+F(B_1+jC),\qquad
 j_{\rm old}=1+\frac{F-1}{C}B_1+Fj\ge1
\]
parametrizes precisely those old witnesses; the quotient is an integer.
It introduces no new variable.

Use the polynomial margin with its known factor 8 removed:
\[
 t=bB^{116},\qquad
 m=K^2(4t-g)-16t(C-KY)^2.
\]
Then \(8m\) is exactly the original Jones margin.  Its strict sign is
unchanged.  The imported ratio interface is now
\[
 a_0\text{ square},\quad a_1\text{ square},\quad J\mid Q,\quad m>0.
\]
Soundness gives \(b\) a power of two, \(8g<G_{\rm cap}\), and the
central-binomial divisibility.  Completeness uses \(32g<G_{\rm cap}\).
All of its numerical hypotheses were established before this invocation.

The strict positive majorants
\[
 V_0=PK,\qquad T_*=k_0(8D^2-4D-1),\qquad V_1=T_*GH
\]
satisfy
\begin{align*}
 V_0^2-a_0&=K^2-1>0,\\
 T_*^2-64DF&=(8D+1)k_0^2-64D>0,\\
 \Delta_1:=V_1^2-64a_1
 &=\bigl((8D+1)k_0^2-64D\bigr)(GH)^2
   +64DF(H^2-1)>0.
\end{align*}
Also \(0<Q/J<Q<3C^2<3D<T_*\le V_1\).  The promised equation is
\begin{align*}
 \Omega&=Jm,\qquad
 \Theta=Jn+Q-\Omega(V_0+V_1).
\end{align*}
\[
 \boxed{F_9=\bigl[\Theta^2-\Omega^2(a_0+64a_1)\bigr]^2
              -256\Omega^4a_0a_1=0.}
\]
If \(m=0\), then \(\Theta=Jn+Q>0\), so there is no zero.  Otherwise
the real norm factors give roots \(r_0^2=a_0\), \(r_1^2=64a_1\) and
\[
 n+Q/J=m(V_0+V_1+r_0+r_1).
\]
The majorants make the parenthesis positive before any integrality claim.
Thus \(m>0\).  The margin gives \(C>K\), and the elementary inequalities
\(A^2>P\), \(D>C^2\), \(F\ge A^2\), \(G\ge A\), \(I\ge G^2\)
give \(a_1>a_0\).  Hence the rational signed sum
\(x=\Theta/\Omega=r_0+r_1\) is nonzero.  From
\(2xr_0=x^2+a_0-64a_1\), both roots are rational and therefore integral;
\(r_1/8\) is integral as well.  The factor equation then yields \(J\mid Q\).
Conversely choose positive roots and set
\(n=m(V_0+V_1+r_0+r_1)-Q/J>0\).  This proves the finalizer equivalence.

\subsection{Exact degree and implementation certificate}
The diagonal degrees are \(r=\deg R=9118\), \(t=\deg N=2403\).
The degree-certificate circuit resolves the near-square cancellation by
\begin{align*}
 &\Theta^2-\Omega^2(a_0+64a_1)\\
 &\quad=E_0^2-2E_0\Omega(V_0+V_1)
 +\Omega^2\bigl[(K^2-1)+\Delta_1+2V_0V_1\bigr],
 \qquad E_0=Jn+Q.
\end{align*}
The unique dominant term is \(\Omega^2\Delta_1\), of degree 1,338,274.
Its square dominates the final subtracted norm term, of degree 1,823,794.
Thus
\[
 \deg F_9=188r+400t+8W+236=2676548,
 \qquad \operatorname{lc}_{\rm diag}(F_9)
 =2^{542}a^{1337768}c^{400}>0.
\]
The diagonal attains the full multivariate upper bound, proving equality.
The positive coordinates replace \(\phi,j,n\) by their positive inputs
minus 1; the natural version instead increments the other six inputs.
These are domain bijections, not integer-square parametrizations.

The supplied checker authenticates both separate whole tuples, rebuilds all
16 primary/comparison and certificate circuits, proves 24 source-evaluator
polynomial identities, and replays 480 actual-tuple modular assignments.
It additionally exercises 43,945 binary-block cases, 540 threshold identities,
Pell-majorant and CRT identities, and finalizer completeness.  These finite
checks audit the implementation; they do not replace the universal proof.
The source and full theorem are detailed in \texttt{nine/PROOF.md}.
